\section{引言}

在格点上抽取部分子分布函数（PDF）$f(x)$ \cite{Feynman:1973xc}
是一个引人关注且颇具挑战性的问题。
在格点上处理 PDF 的通常方法是计算其矩。
然而最近，X.~Ji \cite{Ji:2013dva} 提出了一种
可以把 PDF 作为 $x$ 的函数来计算的方法。

由于 PDF 与光锥 $z^2=0$ 上双定域算符的矩阵元相联系，
这曾是阻碍人们在欧氏空间表述的
格点规范理论中直接计算这些函数的绊脚石。

为克服这一困难，X.~Ji 建议采用纯空间型的间隔 $z=(0,0,0,z_3)$。
此时的函数是准部分子分布（quasi-PDF）$Q(y,p_3)$，它描述强子动量
$p_3$ 分量的分布。
关键在于准部分子分布 $Q(y,p_3)$ 在 \mbox{$p_3 \to \infty$} 极限下
趋于通常的 PDF $f(y)$。
同样的方法也适用于分布振幅（DA）。
格点上准部分子分布的计算结果已见于文献
\cite{Lin:2014zya,Chen:2016utp,Alexandrou:2015rja}，
π 介子准分布振幅的结果见文献 \cite{Zhang:2017bzy}。

本文作者之一（A.R.）的近期论文 \cite{Radyushkin:2016hsy,Radyushkin:2017gjd}
研究了准部分子分布与准分布振幅的非微扰 $p_3$ 演化。该研究基于
虚度分布函数（virtuality distribution functions）的框架
\cite{Radyushkin:2014vla,Radyushkin:2015gpa}。
文献 \cite{Radyushkin:2016hsy,Radyushkin:2017gjd} 中发展的方法建立了
准部分子分布与
“直链接”（straight-link）横动量依赖分布（TMD）
${\cal F} (x, k_\perp^2)$ 之间的联系。
从 TMD 的简单模型出发，可以构造
准部分子分布非微扰演化的模型。所得曲线与格点模拟给出的
\mbox{$p_3$ 演化}图样在定性上一致。

文献 \cite{Radyushkin:2017cyf} 进一步研究了准部分子分布的结构。
该文指出，
当强子运动时，部分子动量 $k_3$
可视为来自两个源头。强子的整体运动贡献
$xp_3$ 部分，它由 TMD ${\cal F} (x, \kappa^2)$ 对第一个宗量即 $x$ 的依赖所支配。
剩余部分 $k_3-xp_3$
则由 TMD 对第二个宗量 $\kappa^2$ 的依赖控制，
后者决定了原始静止系动量分布的形状。
由于准部分子分布具有卷积结构，其 \mbox{$p_3$ 演化}模式相当复杂，
因此必须取较大的 $p_3 \gtrsim 3$ GeV 才能可靠地逼近 PDF 极限。

为了加速收敛，文献 \cite{Radyushkin:2017cyf} 提出了
另一种从格点计算抽取 PDF 的方法，
它基于{\it 伪部分子分布}（pseudo-PDF）$ {\cal P} (x, z_3^2) $的概念。
伪部分子分布把光锥 PDF $f(x)$ 推广到类空间隔
$z=(0,0,0,z_3)$ 情形。它们是
\mbox{{\it Ioffe 时间}\ \cite{Ioffe:1969kf}} \mbox{{\it 分布}\ \cite{Braun:1994jq}}
${\cal M} (\nu, z_3^2)$ 的 Fourier 变换；
后者一般由矩阵元 $  \langle p |   \phi(0) \phi (z)|p \rangle $
作为 $\nu = p_3 z_3$ 和 $z_3^2$ 的函数给出。
与准部分子分布不同，伪部分子分布在一切 $z_3^2$ 取值下都具有“标准的”
$-1 \leq x  \leq 1$ 支撑区。
在 $z_3\to 0$ 极限下它们趋于 PDF，并且在此极限下
呈现以 $1/z_3$ 为演化参数的典型微扰演化。


{如文献~\cite{Radyushkin:2016hsy,Radyushkin:2017gjd} 所讨论的，
伪部分子分布 $ {\cal P} (x, z_3^2) $ 或 Ioffe 时间分布
${\cal M} (\nu, z_3^2)$ 随 $z_3^2$ 的快速非微压下降，
正是准部分子分布 $Q(y, p_3)$ 延迟趋近 PDF $f(y)$ 的原因。}
一个重要的观察是：只要用一个满足 $D(0)=1$、
且其 $z_3^2$ 依赖（平均而言）接近 ${\cal M} (\nu, z_3^2)$ 的适当因子
$D(z_3^2)$ 去除 Ioffe 时间分布 ${\cal M} (\nu, z_3^2)$，
就可以显著削弱其 $z_3^2$ 依赖。
该因子不含 \mbox{$\nu$ 依赖}并满足归一化 $D(0)=1$，
这就保证了比值 ${\cal M} (\nu, z_3^2)/D(z_3^2)$ 在 $z_3^2 \to 0$ 极限下
给出的 PDF 与原函数 ${\cal M} (\nu, z_3^2)$ 在同一极限下给出的相同。


文献 \cite{Radyushkin:2017cyf} 所主张的 $D(z_3^2)$ 的取法，
是令其等于静止系函数 $ {\cal M} (0, z_3^2)$。
这一选择还有一层好处：
${\cal M} (\nu, z_3^2)$ 与 ${\cal M} (0, z_3^2)$ 都含有由规范链接重正化生成的同一乘性
因子 $Z(z_3^2)$，
它在比值中应当消去。

本工作的目标，是按照 \mbox{文献~\cite{Radyushkin:2017cyf}} 中概述的策略，
对质子的 $u$-$d$ PDF 进行探索性的格点计算。
为使本文自洽，我们在第 2、3 节复述
\mbox{文献~\cite{Radyushkin:2017cyf}} 的主要思想；
约化 Ioffe 时间分布的格点抽取方法在第 4 节给出；
数据分析及其解释在第 5 节讨论；
全文总结见第 6 节。
